\documentclass[a4paper,12pt]{book}
\usepackage[top=3cm,left=3cm,right=2cm,bottom=2cm]{geometry} % define margens
\usepackage{setspace} % espaçamento entre linhas
\usepackage[brazil]{babel}
\usepackage[latin1]{inputenc}
\usepackage{amsthm,amsfonts,amssymb,amsmath}
\usepackage{multicol}
\usepackage{wrapfig}
% pacotes gráficos
\usepackage{color}
%\usepackage[dvips]{graphicx}
\usepackage{graphicx,color}
\usepackage{subfigure}
\usepackage{here}
%\usepackage{curves}\curvewarnfalse
%\usepackage{xypic}
\usepackage[colorlinks,plainpages,backref]{hyperref}
%\usepackage{pstricks-add}
%\usepackage{pst-node}
%\usepackage{pst-plot}
%\usepackage[active]{srcltx}
%\usepackage{wrapfig}
%\usepackage{graphicx,graphpap}
\usepackage{array}
\usepackage{fancyhdr}

\usepackage{gfsartemisia-euler}
\usepackage[T1]{fontenc}


\newcommand{\np}{\noindent}
\def\p{\hspace{1.5em}}
\newcommand{\ds}{\displaystyle}

\newtheorem{teo}{Teorema}%[section]
\newtheorem{lema}[teo]{Lema}
\newtheorem{cor}[teo]{Corolário}
\newtheorem{prop}{Proposição}%[section]
\newtheorem{defi}{Definição}%[section]
\newenvironment{dem}{\par \vspace{10pt} \np {\it Prova: }}{$\hfill \square$\vspace{10pt}\par} % demonstra��o
\newenvironment{obs}{\par \vspace{10pt} \scriptsize \np Obs.: }{\vspace{10pt}\par} %observa��o

\newcommand{\N}{\mathbb{N}} % naturais
\newcommand{\Z}{\mathbb{Z}} % inteiros
\newcommand{\R}{\mathbb{R}} % reais
\newcommand{\sen}{\mathrm{sen\,}} % seno, em português
\newcommand{\arcsen}{\mathrm{arcsen\,}} % arco-seno, em português
\newcommand{\Dom}{\mathrm{Dom\,}} % domínio
\newcommand{\Img}{\mathrm{Img\,}} % domínio
\newcommand{\tg}{\mathrm{tg\,}} % seno, em português
\newcommand{\senh}{\mathrm{senh\,}} % seno, em português
\newcommand{\adj}{\mathrm{adj\,}} % adjunta

\begin{document}

\title{Apostila - Pré-Cálculo}
\author{Éden Santana Campos Amorim}
\pagestyle{fancy}
\lhead{Pré-Cálculo}
\chead{Funções básicas}
\rhead{CEFET-MG}

\section*{Funções básicas}

%\begin{multicols}{2}

Listamos a seguir as funções básicas e seus gráficos, ou seja, aquelas dadas por expressões simples, diretas. Dentre elas estão potências, trigonométricas, exponenciais e logaritmos. São apresentadas como na tabela modelo abaixo.\\

%\begin{wraptable}{t}{0.3 \linewidth}
\begin{center}
\begin{tabular}{|c|}
\hline \emph{nome da função} \\ \hline
\emph{gráfico} \\
 \includegraphics[scale=0.2]{geral} \\ \hline
\emph{equação} \\ \hline
\emph{nome da curva} \\ \hline
\end{tabular}
\end{center}
%\end{wraptable}

%\end{multicols}

\subsection*{Funções algébricas}

\begin{center}
\begin{tabular}{|c|c|c|c|}
\hline
Função linear & Função quadrática & Função cúbica & Função recíproca  \\ \hline
\includegraphics[scale=1.1]{linear} & \includegraphics[scale=1.1]{quadratica} & \includegraphics[scale=1.1]{cubica.png} & \includegraphics[scale=0.35]{hiperbole.png} \\ \hline
$y=x$ & $y=x^2$ & $y=x^3$ & $\displaystyle y=\frac{1}{x}=x^{-1}$ \\ \hline
reta & parábola & cúbica & hipérbole \\ \hline
\end{tabular}
\end{center}

\begin{center}
\begin{tabular}{|c|c|c|}
\hline
 Função raiz quadrada & Função raiz cúbica & -  \\ \hline
\includegraphics[scale=1.15]{raiz.png} & \includegraphics[scale=1.05]{raizcubica.png} & \includegraphics[scale=1.5]{semicirc} \\ \hline
$y=\sqrt{x} = x^{\frac{1}{2}}$ & $y=\sqrt[3]{x} = x^{\frac{1}{3}}$ & $\displaystyle y=\sqrt{1-x^2}$ \\ \hline
semi-parábola & cúbica & semi-circunfêrencia \\ \hline
\end{tabular}
\end{center}
\newpage

\subsection*{Funções trigonométricas}

\begin{center}
\begin{tabular}{|c|c|}
\hline
Função seno & Função cosseno \\ \hline
\includegraphics[scale=0.8]{seno} &  \includegraphics[scale=0.8]{cosseno} \\ \hline
$y=\sen x$ & $y=\cos x = \sen(\frac{\pi}{2}-x)$ \\ \hline
senoide & senoide \\ \hline
\end{tabular}
\end{center}

\begin{center}
\begin{tabular}{|c|c|}
\hline
Função tangente & Função secante \\ \hline
 \includegraphics[scale=0.7]{tang} & \includegraphics[scale=0.68]{secante} \\ \hline
 $\displaystyle y=\tan x = \frac{\sen x}{\cos x}$ & $\displaystyle y=\sec x = \frac{1}{\cos x}$ \\ \hline
 - & - \\ \hline
\end{tabular}
\end{center}

\subsection*{Funções exponencial e logaritmo}

\begin{center}
\begin{tabular}{|c|c|}
\hline
Função exponencial & Função logaritmo \\ \hline
\includegraphics[scale=0.7]{exp} & \includegraphics[scale=0.7]{ln} \\ \hline
$y=e^x$ & $y=\ln x$ \\ \hline
- & - \\ \hline
\end{tabular}
\end{center}

\subsection*{Transformações afins}

Abaixo, listamos algumas transformações nas variáveis $x$ e $y$ e como elas afetam o gráfico de um função $y=f(x)$.

\begin{center}
\begin{tabular}{|l|l|} \hline
$f(x) \mapsto f(-x)$ & Reflexão em torno do eixo $y$ \\ \hline
$f(x) \mapsto -f(x)$ & Reflexão em torno do eixo $x$ \\ \hline
$f(x) \mapsto f(x+a)$ & Translação horizontal  \\ 
& direita se $a < 0$, esquerda se $a > 0$ \\ \hline
$f(x) \mapsto f(x)+a$ & Translação vertical, \\
& abaixo se $a < 0$, acima se $a > 0$ \\ \hline
$f(x) \mapsto f(ax)$ & Mudança de escala horizontal \\
& contração se $a > 1$, expansão se $0<a < 1$ \\ \hline
$f(x) \mapsto af(x)$ & Mudança de escala vertical, $a>0$ \\
& expansão se $a >1$, contração se $0<a < 1$ \\ \hline
\end{tabular}
\end{center}

\begin{center}
\begin{tabular}{|c|c|c|}
\hline
Original & Reflexão em $y$ & Reflexão em $x$  \\ \hline
\includegraphics[scale=0.25]{geral} & \includegraphics[scale=0.25]{reflecty} & \includegraphics[scale=0.25]{reflectx} \\ \hline
$y=f(x)$ & $y=f(-x)$ & $y=-f(x)$ \\ \hline
\end{tabular}
\end{center}

\begin{center}
\begin{tabular}{|c|c|c|c|}
\hline
Translação direita & Translação esquerda & Translação acima & Translação abaixo \\ \hline
\includegraphics[scale=0.25]{transh} & \includegraphics[scale=0.25]{transh2} & \includegraphics[scale=0.25]{transv} & \includegraphics[scale=0.25]{transv2}\\ \hline
$y=f(x+a), a< 0$ & $y=f(x+a), a> 0$ & $y=f(x)+a, a>0$ & $y=f(x)+a, a<0$ \\ \hline
\end{tabular}
\end{center}

\begin{center}
\begin{tabular}{|c|c|c|c|}
\hline
Expansão horizontal & Contração horizontal & Expansão vertical & Contração vertical \\ \hline
\includegraphics[scale=0.25]{scaleh} & \includegraphics[scale=0.25]{scaleh2} & \includegraphics[scale=0.25]{scalev} & \includegraphics[scale=0.25]{scalev2} \\ \hline
$y=f(ax), 0<a< 1$ & $y=f(ax), a> 1$ & $y=af(x), a>1$ & $y=af(x), 0<a<1$  \\ \hline
\end{tabular}
\end{center}

{\scriptsize
\np \emph{Observações:}

\np $\bullet$ Toda função quadrática pode ser obtida a partir de $y=x^2$ por transformações afins usando o completamento de quadrados: $x^2+2ax+a^2=(x+a)^2$.

\np $\bullet$ Assim como a função cosseno pode ser obtida a partir de seno por transformações afins, as funções cotangente e cossecante também podem ser obtidas a partir da tangente e secante, respectivamente: $\cot x = \tan(\frac{\pi}{2}-x)$ e $\csc x = \sec(\frac{\pi}{2}-x)$.

\np $\bullet$ As funções exponenciais e logaritmos em outras bases também podem ser obtidas a partir das funções de base $e$ por transformações afins, lembrando das fórmulas de mudança de base: $a^x = e^{x \ln a}$ e $\log_a x = \frac{\ln x}{\ln a}$.}


\end{document}

